\section{Conclusions}
\label{sec:conclusions}

We have investigated the tidal origin of stellar disk warps using a tilted-ring dynamical model with 60 concentric rings, evolving under three torques: disk self-gravity coupling, an oblate NFW dark matter halo, and the quadrupole tidal field of an orbiting satellite galaxy. Our main conclusions are as follows.

\begin{enumerate}
    \item \textbf{Impulsive excitation.} Warp growth is impulsive rather than continuous: the tilt amplitude increases in a characteristic staircase pattern, with each step synchronized to a pericentric passage of the satellite (Section~\ref{sec:fiducial}). Between passages, the warp amplitude plateaus. This impulsive nature has implications for interpreting the MW warp as a snapshot of ongoing satellite interaction.
    
    \item \textbf{Radial structure.} The tidal torque scales as $R^2$, producing a monotonically increasing tilt profile that peaks at the disk edge. The inner disk ($R\lesssim4$\,kpc) is stabilized by the bending stiffness of disk self-gravity (Section~\ref{sec:fiducial}).
    
    \item \textbf{Linear mass scaling.} The warp amplitude scales linearly with satellite mass over two orders of magnitude ($M_\mathrm{sat}=10^{10}$--$5\times10^{11}\,\mathrm{M}_\odot$), as predicted by the quadrupole tidal torque formula (Section~\ref{sec:param_mass}).
    
    \item \textbf{$\sin 2i$ inclination dependence.} The warp peaks at orbital inclination $i\approx45^\circ$ and vanishes for coplanar ($i=0^\circ$) or polar ($i=90^\circ$) orbits, following a $\sin 2i$ dependence. Polar orbits produce an oscillatory warp rather than a monotonically growing one (Section~\ref{sec:param_incl}).
    
    \item \textbf{Coherent precession.} The line of nodes precesses coherently at ${\sim}7\,\mathrm{deg\,Gyr^{-1}}$, maintained by disk self-gravity coupling that resists the differential winding imposed by the halo. This rate is consistent with \textit{Gaia}-based measurements \citep{poggio2020, cheng2020}.
    
    \item \textbf{Torque hierarchy.} At $R>2$\,kpc, the tidal torque dominates over disk self-gravity by 1--3 orders of magnitude and over the halo torque by ${\sim}3$ orders of magnitude. Consequently, the warp amplitude is insensitive to halo flattening ($q=0.8$--$1.0$) in the presence of a strong tidal perturber (Section~\ref{sec:param_qhalo}).
    
    \item \textbf{Comparison with observations.} The fiducial model (peak tilt $1.70^\circ$, outer-disk tilt $1.34^\circ$) lies at the lower end of the observed MW stellar warp ($2^\circ$--$5^\circ$; \citealt{chen2019, skowron2019}). The deficit is likely attributable to our neglect of the dynamical halo response, disk truncation at $R=18$\,kpc, and the absence of a second perturber (Section~\ref{sec:comparison_mw}).
\end{enumerate}

The tilted-ring framework, while simplified, provides a physically transparent and computationally efficient tool for understanding the parametric dependencies of tidal warp formation. Future work should incorporate a live dark matter halo to quantify the amplification of the tidal response, extend the disk to larger radii, and include gas dynamics to capture the differential stellar/gaseous warp response.
